%% ============================================================
%% Rozdział 1 — Wstęp i cel pracy
%% ============================================================
\chapter{Wstęp i cel pracy}\label{chapter:1_Introduction}
Postępujący rozwój technologii \dots \: W niniejszym rozdziale należy przedstawić problematykę, która będzie analizowana lub rozwiązywana w ramach pracy dyplomowej, wraz z uzasadnieniem podjęcia tematu. Powinno się tu syntetycznie omówić dotychczasowy stan wiedzy oraz aktualne rozwiązania odnoszące się do wybranego zagadnienia.

Zaleca się ujęcie teoretycznego opisu problemu, definicji podstawowych pojęć, a także przeglądu istniejących rozwiązań – z uwzględnieniem ich mocnych i słabych stron. Istotne jest również jednoznaczne sformułowanie tematu pracy, zgodnego z zaakceptowaną kartą tematu w systemie mojaPG. Całość powinna być przygotowana z dbałością o estetykę redakcyjną, poprawność językową i terminologiczną, a także z zachowaniem spójności merytorycznej.

\section{Przegląd literatury}
Od wielu lat obserwuje się dynamiczny wzrost liczby publikacji poświęconych \dots \: W niniejszym rozdziale należy przygotować syntetyczny przegląd literatury obejmujący od 10 do 20 pozycji. Powinien on uwzględniać zarówno klasyczne opracowania stanowiące fundament danej dziedziny, jak i najnowsze publikacje z ostatnich pięciu lat. Celem tego przeglądu jest krótkie, lecz spójne przedstawienie dotychczasowej wiedzy oraz aktualnych rozwiązań odnoszących się do problematyki pracy, a także ich wstępna analiza i ocena.

Cytowania w tekście należy wprowadzać za pomocą polecenia \verb|\cite{klucz}|, np. \cite{DeanGhemawat2004MapReduce}. Bibliografia powinna być przygotowana zgodnie z wytycznymi normy \textbf{PN-ISO 690:2012} , z wykorzystaniem pliku \verb|_bibliography.bib|, w którym definiuje się wszystkie źródła dołączone do pracy. Pozycje literatury \cite{HorowitzHill2015AoE} zostaną zebrane w \hyperref[page:Bibliography]{Wykazie literatury}, zgodnie z wymaganiami redakcyjnymi i zasadami poprawnego przygotowania pracy dyplomowej.

Numeracja bibliografii jest automatyczna względem kolejności występowania w tekście pracy \cite{Kalman1960NewApproach}. Jeśli konieczna jest inna forma numeracji, można ją zmodyfikować w pakiecie \texttt{biber}, np. w pliku \verb|_formatting.cls|, choć zaleca się korzystanie z domyślnego ustawienia.

\section{Cel i zakres pracy}
Tematem niniejszej pracy jest \dots \: W tym miejscu należy jednoznacznie i zwięźle sformułować \textbf{temat} oraz precyzyjny cel pracy, a w razie potrzeby podać również tezę (hipotezę). Następnie należy syntetycznie przedstawić dotychczasowe dokonania w danej dziedzinie, przyjęte założenia techniczne oraz krótko zapowiedzieć zawartość kolejnych rozdziałów.

We wprowadzeniu warto zarysować aktualny stan wiedzy i stosowane rozwiązania techniczne, wskazać istniejące luki oraz kluczowe wyzwania projektowo-badawcze. Należy także uzasadnić, że dalsze części pracy rozwijają temat w sposób oryginalny i zasadny.

W przypadku prac wieloautorskich konieczne jest jednoznaczne przypisanie autorstwa do poszczególnych części. Każdy podrozdział powinien być opracowany przez jednego autora, a autorów rozdziałów lub podrozdziałów należy wymienić w odpowiednim miejscu dokumentu.
